\documentclass[a4paper,11pt]{article}
\usepackage[T2A]{fontenc}
\usepackage[utf8]{inputenc}
\usepackage[russian]{babel}
\usepackage{amsmath,amssymb,mathtools}
\usepackage[margin=20mm,headheight=14pt,footskip=13mm]{geometry}
\usepackage{microtype}
\usepackage{booktabs,array,tabularx}
\usepackage{xcolor}
\usepackage{tikz}
\usepackage{fancyhdr}
\usepackage[hidelinks]{hyperref}
\definecolor{accent}{HTML}{147C83}
\definecolor{ink}{HTML}{243238}
\definecolor{signal}{HTML}{AE3533}
\definecolor{idea}{HTML}{205E9B}
\setlength{\parindent}{0pt}
\setlength{\parskip}{6pt}
\setlength{\emergencystretch}{2em}
\linespread{1.09}
\pagestyle{fancy}
\fancyhf{}
\fancyfoot[C]{\footnotesize Лёвин Артём Александрович эксклюзивно для Николь Саркисьянц}
\fancyfoot[R]{\footnotesize\thepage}
\renewcommand{\headrulewidth}{0pt}
\renewcommand{\footrulewidth}{0.3pt}
\newcommand{\errorlabel}{\textcolor{signal}{\textbf{Ошибка:}}}
\newcommand{\idealabel}{\textcolor{idea}{\textbf{Ключевая идея:}}}
\newcommand{\resultlabel}{\textcolor{accent}{\textbf{Итог:}}}
\newcommand{\sinx}{\sin 2x}
\begin{document}
\thispagestyle{empty}
\vspace*{43mm}
\begin{center}
{\Large\color{accent} МАТЕМАТИКА}\par
\vspace{17mm}
{\Huge\bfseries Ревью домашней работы}\par
\vspace{10mm}
{\large Тригонометрические преобразования и уравнения}\par
\vspace{25mm}
{\large 26 октября 2026 г.}\par
\vspace{15mm}
{\large Лёвин Артём Александрович}\par
\end{center}
\vfill
\begin{center}
{\small Лёвин Артём Александрович эксклюзивно для Николь Саркисьянц}\par
\vspace{10mm}
\begin{tikzpicture}
\draw[accent!55,line width=1.2pt] (-4,0) .. controls (-2,0.6) and (-0.9,-0.4) .. (0,0)
 .. controls (1,0.4) and (2.1,-0.6) .. (4,0.05);
\end{tikzpicture}
\end{center}
\newpage
\section*{Сводная таблица}
\smallskip
В таблице указаны только задания, в которых обнаружено нарушение условия или математической записи. По номеру можно перейти к подробному разбору.
\par\vspace{4mm}
\begingroup
\renewcommand{\arraystretch}{1.35}
\small
\begin{tabularx}{\linewidth}{@{}p{0.85cm}p{2.25cm}X p{2.40cm}p{2.85cm}@{}}
\toprule
\textbf{№} & \textbf{Тема} & \textbf{Суть ошибки} & \textbf{Ваш ответ} & \textbf{Правильный результат}\\
\midrule
\hyperref[task:9]{\textbf{9}} & Группировка & К выражению без основания добавлено равенство нулю. & Произведение \(=0\) & Тождество: исходное выражение \(=\) произведению \\
\addlinespace
\hyperref[task:10]{\textbf{10}} & Уравнение & Не представлен второй из двух требуемых способов решения. & Корни найдены одним способом & Те же корни; два полных способа \\
\bottomrule
\end{tabularx}
\endgroup
\vfill
\begin{center}
\textcolor{accent}{\rule{\linewidth}{0.5pt}}\\[5pt]
\small Разбирайте сначала точку, в которой возникло затруднение, затем повторяйте полное решение.
\end{center}
\newpage
\thispagestyle{empty}
\vspace*{0.34\textheight}
\begin{center}
{\Huge\bfseries\color{accent} Разбор ошибок}\par
\vspace{12mm}
{\large Пошаговые исправления и задания на закрепление}
\end{center}
\newpage
\phantomsection\label{task:9}
\section*{Задание 9. Группировка и равенство выражений}
\textbf{Текст задания.} Разложите на множители:
\[
4\sin^3 2x-2\sin^2 2x-2\sin 2x+1.
\]
\textbf{Ваше решение.} В записи после исходного выражения поставлено \(=0\), затем выполнена группировка:
\begin{align*}
&2\sin^2 2x(2\sin 2x-1)-(2\sin 2x-1)=0,\\
&(2\sin 2x-1)(2\sin^2 2x-1)=0.
\end{align*}
\textbf{Ваш ответ.} \((2\sin 2x-1)(2\sin^2 2x-1)=0\).

\errorlabel{} К выражению приписано \(=0\), хотя в условии требуется только разложение на множители. Сама группировка выполнена верно.

\textbf{Где возникает ошибка.} Последний правильный объект --- данное выражение. Первый неправильный переход --- замена его на уравнение:
\[
4\sin^3 2x-2\sin^2 2x-2\sin 2x+1\ \color{signal}{=0}.
\]
Условие \(=0\) отсутствует в задании. После этого все равенства с нулём относятся к самостоятельно введённому уравнению.

\textbf{Почему это неверно.} Разложить выражение --- значит найти равное ему произведение при любом допустимом \(x\). Равенство произведения нулю утверждает, что выражение имеет нулевое значение. Это уже другое условие и оно вообще говоря неверно. Например, при \(x=0\):
\[
4\sin^3 0-2\sin^2 0-2\sin 0+1=1\ne 0.
\]
\idealabel{} При группировке ставьте знак равенства между \emph{исходным выражением} и его преобразованной формой, не добавляя уравнения.

\textbf{Исправление от последнего правильного шага.} Положим \(t=\sin 2x\). Группируем слагаемые:
\begin{align*}
4t^3-2t^2-2t+1
&=(4t^3-2t^2)+(-2t+1)\\
&=2t^2(2t-1)-(2t-1)\\
&=(2t-1)(2t^2-1).
\end{align*}
\textbf{Полное правильное решение.} Возвращаемся к \(t=\sin 2x\):
\[
\boxed{\begin{aligned}
&4\sin^3 2x-2\sin^2 2x-2\sin 2x+1\\
&\qquad=(2\sin 2x-1)(2\sin^2 2x-1).
\end{aligned}}
\]
\textbf{Проверка.} При раскрытии скобок получаем
\[
4\sin^3 2x-2\sin 2x-2\sin^2 2x+1,
\]
что совпадает с исходным выражением. Кроме того, при \(x=0\) обе части тождества равны единице.

\resultlabel{} Множители найдены верно; в ответе нужно записать тождество без условия \(=0\).

\textbf{Задача на закрепление.} Разложите на множители выражение
\[
6\sin^3 x-3\sin^2 x-2\sin x+1.
\]
\newpage
\phantomsection\label{task:10}
\section*{Задание 10. Два способа решения}
\textbf{Текст задания.} Решите \emph{двумя способами}: через формулы двойного угла и через переход от произведения к сумме:
\[
\cos 2x\,\sin 4x-\sin\frac{5\pi}{6}\cos 4x
=\cos\left(2x-\frac{\pi}{2}\right).
\]
\textbf{Ваше решение.} Вы верно нашли
\[
\sin\frac{5\pi}{6}=\frac12,\qquad
\cos\left(2x-\frac{\pi}{2}\right)=\sin 2x,
\]
затем применили формулы двойного угла, привели уравнение к
\[
(2\sin 2x-1)(2\sin^2 2x-1)=0
\]
и выписали шесть семейств корней: два для \(\sin 2x=\frac12\) и четыре для \(\sin 2x=\pm\frac{\sqrt2}{2}\). Все эти семейства верны.

\textbf{Ваш ответ.} Корни найдены, но приведён только первый способ.

\errorlabel{} Не выполнено обязательное требование решить уравнение двумя различными способами. Вычислительной ошибки в представленном способе нет.

\textbf{Где именно возникает недочёт.} Последний правильный шаг --- нахождение всех шести семейств корней. После него работа завершается, хотя условие прямо требует ещё одного решения через формулу преобразования произведения в сумму.

\textbf{Почему это важно.} Два способа здесь показывают связь разных тригонометрических формул. Они должны приводить к одному и тому же множеству корней. При оформлении необходимо предъявить оба преобразования независимо.

\idealabel{} При переходе от произведения к сумме исходное уравнение сразу приводится к произведению двух множителей. Это более короткий путь и одновременно проверка первого.

\subsection*{Способ 1. Через двойной угол}
Сначала используем уже установленные значения:
\[
\cos 2x\,\sin 4x-\frac12\cos 4x=\sin 2x.
\]
Обозначим \(t=\sin 2x\). Тогда
\[
\sin 4x=2\sin 2x\cos 2x,\quad
\cos^2 2x=1-t^2,\quad
\cos 4x=1-2t^2.
\]
Подставляем формулы, не пропуская слагаемых:
\begin{align*}
2t(1-t^2)-\tfrac12(1-2t^2)&=t,\\
2t-2t^3-\tfrac12+t^2-t&=0,\\
-2t^3+t^2+t-\tfrac12&=0.
\end{align*}
Умножим обе части на \(-2\) и сгруппируем:
\begin{align*}
4t^3-2t^2-2t+1&=0,\\
2t^2(2t-1)-(2t-1)&=0,\\
(2t-1)(2t^2-1)&=0.
\end{align*}
Произведение равно нулю, когда хотя бы один множитель равен нулю:
\[
\sin 2x=\frac12
\quad\text{или}\quad
2\sin^2 2x-1=0.
\]
Во втором случае
\[
2\sin^2 2x-1=-\cos 4x,
\]
поэтому его корни задаются условием \(\cos 4x=0\). Решаем две простые серии уравнений:
\begin{align*}
\sin 2x=\tfrac12
&\iff
\begin{cases}
2x=\frac\pi6+2\pi n,\\
2x=\frac{5\pi}6+2\pi n,
\end{cases}\\[-2pt]
&\iff
\begin{cases}
x=\frac\pi{12}+\pi n,\\x=\frac{5\pi}{12}+\pi n;
\end{cases}\\[3pt]
\cos 4x=0
&\iff 4x=\frac\pi2+\pi k\\
&\iff x=\frac\pi8+\frac{\pi k}4,
\end{align*}
где \(n,k\in\mathbb Z\).

\subsection*{Способ 2. Через переход от произведения к сумме}
Начнём с исходного уравнения после вычисления постоянных тригонометрических значений:
\[
\cos 2x\,\sin 4x-\frac12\cos 4x=\sin 2x.
\]
Используем формулу
\[
\sin\alpha\cos\beta
=\frac12\bigl(\sin(\alpha+\beta)+\sin(\alpha-\beta)\bigr).
\]
При \(\alpha=4x\), \(\beta=2x\) получаем
\[
\sin 4x\cos 2x=\frac12(\sin 6x+\sin 2x).
\]
Теперь преобразуем уравнение:
\begin{align*}
\frac12(\sin 6x+\sin 2x)-\frac12\cos 4x&=\sin 2x,\\
\sin 6x-\sin 2x-\cos 4x&=0.
\end{align*}
Для разности синусов применим формулу
\[
\sin A-\sin B
=2\sin\frac{A-B}{2}\cos\frac{A+B}{2}.
\]
Подставляя \(A=6x\), \(B=2x\), получаем
\[
\sin 6x-\sin 2x=2\sin 2x\cos 4x.
\]
Следовательно,
\begin{align*}
2\sin 2x\cos 4x-\cos 4x&=0,\\
\boxed{\cos 4x\,(2\sin 2x-1)}&\boxed{=0}.
\end{align*}
Решаем \(\cos 4x=0\) или \(\sin 2x=\tfrac12\). Получаем в точности те же семейства корней, что в первом способе:
\[
\boxed{\begin{gathered}
x=\frac\pi8+\frac{\pi k}4,\quad k\in\mathbb Z,\\[4pt]
x=\frac\pi{12}+\pi n\quad\text{или}\quad
x=\frac{5\pi}{12}+\pi n,\quad n\in\mathbb Z.
\end{gathered}}
\]
\textbf{Проверка.} Все переходы во втором способе являются тождественными; деления на выражение, зависящее от \(x\), нет. Поэтому найденные серии не теряются и посторонних корней не возникает. Для дополнительной проверки подставим \(x=\frac\pi8\):
\[
\cos\frac\pi4\sin\frac\pi2-\frac12\cos\frac\pi2
=\frac{\sqrt2}{2}
=\sin\frac\pi4.
\]
Четыре отдельные серии, выписанные Вами для \(\sin 2x=\pm\frac{\sqrt2}2\), объединяются краткой записью \(x=\frac\pi8+\frac{\pi k}4\).

\resultlabel{} Оба способа дают один и тот же полный набор корней. При задании с формулировкой «двумя способами» нужно оформить каждый из них.

\textbf{Задача на закрепление.} Решите двумя способами --- через формулы двойного угла и через переход от произведения к сумме:
\[
\cos 3x\,\sin 6x-\sin\frac\pi6\cos 6x
=\cos\left(3x-\frac\pi2\right).
\]
\end{document}
